% ============================================================================
% Fichier  : partie4/P04_C15_forensique_nuage.tex : FORN (LPI342)
% Sources : ISO/IEC 27037, 27042 ; NIST (cloud). Opération TSANAGA.
% ============================================================================

\chapter{Forensique du Nuage et des Applications}
\label{chap:nuage}
\chaptermark{Forensique du nuage}

\epigraph{%
    « Dans le nuage, la preuve n'est pas chez vous~;\\
    il faut la demander, dans les formes. »%
}{Principe du cours}

\niveauChapitre{Expert}{Chapitre~\ref{chap:mobile}~; notions de services en
ligne.}

\begin{boxobjectifs}
\begin{itemize}
    \item Comprendre les défis du forensique dans le nuage.
    \item Identifier les traces côté client et côté service.
    \item Obtenir une preuve auprès d'un hébergeur, dans les formes.
    \item Analyser les artefacts d'applications (réseaux sociaux, messageries).
    \item Respecter le droit et la coopération internationale.
\end{itemize}
\end{boxobjectifs}

\bigskip

% ============================================================================
\section*{Question directrice : le fichier a été envoyé vers un service d'hébergement. Comment en obtenir la preuve ?}
% ============================================================================

L'envoi du 5~janvier (Chapitre~\ref{chap:reseau}) pointe vers un service
d'hébergement personnel, hébergé à l'étranger. La preuve n'est plus sur un
support saisi, mais chez un tiers. Ce chapitre traite de ce cas.

\medskip
\textbf{Plan du chapitre.}
\begin{enumerate}
    \item Les défis du nuage (section~\ref{sec:c15-defis}).
    \item Côté client, côté service (section~\ref{sec:c15-cotes}).
    \item Obtenir dans les formes (section~\ref{sec:c15-obtenir}).
    \item Les applications (section~\ref{sec:c15-apps}).
\end{enumerate}

\separateur

% ============================================================================
\section{Les Défis du Nuage}
\label{sec:c15-defis}
% ============================================================================

Dans le nuage, la donnée est répartie, chez un tiers, souvent à l'étranger.
L'enquêteur ne peut ni la saisir physiquement, ni l'imager comme un disque. Il
doit la \emph{demander} au service, dans les formes légales, et composer avec
la localisation des données et la coopération internationale.

% ============================================================================
\section{Côté Client, Côté Service}
\label{sec:c15-cotes}
% ============================================================================

\begin{tableaucours}{L{3.6cm}L{8.6cm}}{%
    \tableauHeader{Côté} &
    \tableauHeader{Traces accessibles}
}
Client (le poste) & Fichiers de synchronisation, caches, jetons, historiques \\
Service (l'hébergeur) & Comptes, journaux d'accès, fichiers, sauvegardes \\
\end{tableaucours}
\vspace{-0.8em}
\begin{center}{\small\itshape Tableau~15.1. Deux côtés de la preuve en nuage}\end{center}

Souvent, le poste saisi garde déjà des traces du service (dossier de
synchronisation, jeton d'accès), exploitables sans attendre le tiers.

% ============================================================================
\section{Obtenir dans les Formes}
\label{sec:c15-obtenir}
% ============================================================================

\begin{boxjuris}
Obtenir une preuve auprès d'un hébergeur étranger passe par les voies
légales~: réquisition, coopération judiciaire internationale (Convention de
Budapest), demande au point de contact du service. Accéder directement au
compte, même avec un jeton trouvé, hors de ce cadre, exposerait à
l'irrecevabilité et à l'illégalité.
\end{boxjuris}

\begin{boxalert}
\textbf{Violation typique.} Se connecter au service avec un jeton d'accès
trouvé sur le poste, pour « aller voir ».

\smallskip
\textbf{Règle~:} un jeton trouvé établit un lien~; il ne vaut pas autorisation
d'accéder au service. On documente le jeton comme trace locale, et l'on
obtient le contenu du service par la voie légale.
\end{boxalert}

% ============================================================================
\section{Les Applications}
\label{sec:c15-apps}
% ============================================================================

\begin{boxscenario}{Opération TSANAGA. La trace de synchronisation}
\textbf{Constatation.} Sur le poste saisi, le dossier de synchronisation du
service d'hébergement contient une copie locale du fichier clients et un
journal de téléversement daté du 5~janvier.

\medskip
\textbf{Analyse} (confiance élevée). La trace locale, recevable (poste saisi,
chaîne continue), établit le téléversement sans attendre l'hébergeur. La
demande au service est engagée en parallèle, dans les formes, pour corroborer.
\end{boxscenario}

\begin{boxmetier}{Analyste forensique et juriste TIC}
\textbf{Chercher d'abord côté client.} Avant d'engager une longue coopération
internationale, l'analyste exploite les traces du service déjà présentes sur
le poste saisi~: elles sont souvent suffisantes et immédiatement recevables.
\end{boxmetier}

\separateur

\begin{boxresume}
\textbf{Ce qu'il faut retenir de ce chapitre~:}
\begin{itemize}
    \item \textbf{Le nuage}~: donnée répartie, chez un tiers, souvent à
          l'étranger.
    \item \textbf{Côté client}, le poste garde souvent des traces du service.
    \item \textbf{Côté service}, on obtient dans les formes (réquisition,
          Budapest).
    \item \textbf{Un jeton trouvé} n'autorise pas l'accès au service.
    \item \textbf{On exploite d'abord les traces locales}, recevables.
\end{itemize}
\end{boxresume}

\begin{boxinfo}
\textbf{Points de contrôle de l'élève.} Je sais~: expliquer les défis du
nuage~; distinguer traces client et service~; obtenir une preuve dans les
formes~; exploiter une trace locale~; respecter le droit.
\end{boxinfo}

\bigskip

\begin{boxexo}[title={\ding{45}\quad Exercice~15.1. Défi \niveaubase}]
Pourquoi ne peut-on pas imager une donnée dans le nuage comme un disque~?
\end{boxexo}

\begin{boxexo}[title={\ding{45}\quad Exercice~15.2. Deux côtés \niveaubase}]
Quelles traces trouve-t-on côté client~? Côté service~?
\end{boxexo}

\begin{boxexo}[title={\ding{45}\quad Exercice~15.3. Jeton \niveaubase}]
Un jeton trouvé autorise-t-il à se connecter au service~?
\end{boxexo}

\begin{boxexo}[title={\ding{45}\quad Exercice~15.4. Dans les formes \niveaumoyen}]
Par quelles voies obtient-on une preuve auprès d'un hébergeur étranger~?
\end{boxexo}

\begin{boxexo}[title={\ding{45}\quad Exercice~15.5. TSANAGA \niveaumoyen}]
Quelle trace locale établit le téléversement, sans attendre l'hébergeur~?
\end{boxexo}

\begin{boxexo}[title={\ding{45}\quad Exercice~15.6. Budapest \niveaumoyen}]
Quel rôle joue la Convention de Budapest ici~?
\end{boxexo}

\begin{boxexo}[title={\ding{45}\quad Exercice~15.7. Synchronisation \niveaumoyen}]
Pourquoi un dossier de synchronisation est-il précieux~?
\end{boxexo}

\begin{boxexo}[title={\ding{45}\quad Exercice~15.8. Priorité \niveaumoyen}]
Pourquoi exploiter d'abord les traces locales~?
\end{boxexo}

\begin{boxexo}[title={\ding{45}\quad Exercice~15.9. Demande \niveauexpert}]
Rédigez la demande de preuve adressée au service d'hébergement.
\end{boxexo}

\begin{boxexo}[title={\ding{45}\quad Exercice~15.10. Bilan domaines \niveauexpert}]
Récapitulez ce que chaque domaine (disque, mémoire, réseau, mobile, nuage) a
apporté à TSANAGA.
\end{boxexo}

\bigskip

\begin{boxinfo}
\textbf{Pour aller plus loin.}
\begin{itemize}
    \item ISO/IEC 27037~\cite{iso27037}, 27042~\cite{iso27042}~;
          Budapest~\cite{budapest2001}.
\end{itemize}
\end{boxinfo}

\begin{boxinfo}
\textbf{Transition.} Un suspect averti efface ses traces~: la Partie~V traite
des défis, à commencer par l'anti-forensique
(Chapitre~\ref{chap:anti}).
\end{boxinfo}
